\chapter{Derivadas e convexidade}\label{ch:g12:deriv}

A \omterm{def:g12:deriv:derivative}{derivada} mede a taxa de variação instantânea de uma \omterm{def:g11:func:function}{função}; é o
\omterm{def:g10:reffunc:affine}{coeficiente angular} da reta \omterm{def:g12:deriv:tangent}{tangente} ao seu \omterm{def:g10:functions:graph}{gráfico}. Este capítulo revê e
estende as regras de derivação, liga o sinal de $f'$ à variação de $f$ e
apresenta a \omterm{def:g12:deriv:convex}{convexidade}, governada pela segunda \omterm{def:g12:deriv:derivative}{derivada}.

\section{A derivada}

\begin{definition}[Derivada em um ponto]\label{def:g12:deriv:derivative}
Seja $f$ definida em um \omterm{def:g10:numbers:interval}{intervalo} $I$ e seja $a \in I$. A \omterm{def:g11:func:function}{função} $f$ é
\emph{derivável em $a$}\index{derivada}\index{derivável} se o quociente de
diferenças
\[
\frac{f(a+h) - f(a)}{h}
\]
tem limite finito quando $h \to 0$. Esse limite chama-se \emph{derivada de
$f$ em $a$} e escreve-se $f'(a)$. Se $f$ é derivável em todo ponto de
$I$, a \omterm{def:g11:func:function}{função} $f' \colon x \mapsto f'(x)$ é a \emph{derivada} de $f$.
\end{definition}

\begin{definition}[Reta tangente]\label{def:g12:deriv:tangent}
Se $f$ é \omterm{def:g12:deriv:derivative}{derivável} em $a$, a \emph{tangente}\index{tangente} à curva de
$f$ no ponto $(a, f(a))$ é a reta de \omterm{def:g10:algebra:equation}{equação}
\[
y = f(a) + f'(a)(x - a).
\]
\end{definition}

\begin{omfigure}
\begin{tikzpicture}
\begin{axis}[omaxis, clip=false,
  width=9.6cm, height=6.8cm,
  xmin=-0.4, xmax=3.4, ymin=-0.4, ymax=4.3,
  xtick={1,2.5}, xticklabels={$a$,$a+h$}, ytick=\empty]
  \addplot[omProp!40, thick, domain=0.2:2.7] {1 + 1.75*(x-1)};
  \addplot[omProp!70, thick, domain=0.07:2.95] {1 + 1.5*(x-1)};
  \addplot[omProp, thick, domain=-0.12:3.3] {1 + 1.25*(x-1)};
  \addplot[omThm, thick, domain=-0.4:3.3] {x};
  \addplot[omDef, thick, domain=-0.4:2.72] {x^2/2 + 0.5};
  \draw[black, dashed, thin] (axis cs:1,1) -- (axis cs:2.5,1)
    node[midway, below, font=\small] {$h$};
  \draw[black, dashed, thin] (axis cs:2.5,1) -- (axis cs:2.5,3.625)
    node[midway, right, font=\small] {$f(a+h)-f(a)$};
  \fill (axis cs:2.5,3.625) circle (1.5pt);
  \fill (axis cs:2,2.5) circle (1.5pt);
  \fill (axis cs:1.5,1.625) circle (1.5pt);
  \fill (axis cs:1,1) circle (1.5pt)
    node[above left, font=\small] {$A$};
  \node[omThm, font=\small, anchor=west] at (axis cs:3.38,3.3) {tangente};
\end{axis}
\end{tikzpicture}

{\small À medida que $h$ diminui, a \omterm{def:g11:deriv:rate}{secante} por $A = (a, f(a))$ e
$(a+h, f(a+h))$ (laranja) gira rumo à \omterm{def:g12:deriv:tangent}{tangente} em $A$ (vermelha): seu
\omterm{def:g10:reffunc:affine}{coeficiente angular} $\frac{f(a+h)-f(a)}{h}$ tende a $f'(a)$.}
\end{omfigure}

\begin{proposition}\label{prop:g12:deriv:diffcont}
Uma \omterm{def:g11:func:function}{função} \omterm{def:g12:deriv:derivative}{derivável} em $a$ é \omterm{def:g12:limcont:continuity}{contínua} em $a$. A recíproca é falsa.
\end{proposition}

\begin{proof}
Para $h \neq 0$ pequeno,
$f(a+h) - f(a) = h \cdot \frac{f(a+h)-f(a)}{h}$. Quando $h \to 0$, o lado
direito tende a $0 \times f'(a) = 0$, logo $f(a+h) \to f(a)$. Para a
recíproca, $x \mapsto \abs{x}$ é \omterm{def:g12:limcont:continuity}{contínua} em $0$, mas seu quociente de
diferenças em $0$ vale $\frac{\abs h}{h} = \pm 1$ conforme o sinal de $h$,
e não tem limite.
\end{proof}

\subsection{Regras de derivação}

\begin{proposition}[Operações]\label{prop:g12:deriv:operations}
Sejam $u, v$ deriváveis em $I$ e $\lambda \in \R$. Então $u + v$,
$\lambda u$ e $uv$ são deriváveis em $I$, assim como $\frac{u}{v}$ onde
$v \neq 0$, e
\[
(u+v)' = u' + v', \quad (\lambda u)' = \lambda u', \quad
(uv)' = u'v + uv', \quad
\left(\frac{u}{v}\right)' = \frac{u'v - uv'}{v^2}.
\]
\end{proposition}

\begin{proof}[Demonstração da regra do produto]
Escreva o quociente de diferenças de $uv$ em $a$ como
\[
\frac{u(a+h)v(a+h) - u(a)v(a)}{h}
= \frac{u(a+h) - u(a)}{h}\,v(a+h) + u(a)\,\frac{v(a+h) - v(a)}{h}.
\]
Quando $h \to 0$, $v(a+h) \to v(a)$ por \omterm{def:g12:limcont:continuity}{continuidade}
(\cref{prop:g12:deriv:diffcont}), de modo que o lado direito tende a
$u'(a)v(a) + u(a)v'(a)$. As outras regras se demonstram de modo análogo.
\end{proof}

\begin{theorem}[Regra da cadeia]\label{thm:g12:deriv:chain}
\index{regra da cadeia}
Seja $u$ \omterm{def:g12:deriv:derivative}{derivável} em $I$ com valores em $J$, e $g$ \omterm{def:g12:deriv:derivative}{derivável} em $J$.
Então $g \circ u$ é \omterm{def:g12:deriv:derivative}{derivável} em $I$ e
\[
(g \circ u)' = (g' \circ u) \cdot u' .
\]
Em particular, para $u$ \omterm{def:g12:deriv:derivative}{derivável}:
\[
(u^n)' = n\,u'\,u^{n-1} \ (n \in \Z), \qquad
\left(\sqrt{u}\right)' = \frac{u'}{2\sqrt{u}} \ (u > 0), \qquad
\left(\eu^{u}\right)' = u'\,\eu^{u}.
\]
\end{theorem}

\begin{proof}[Esboço da demonstração]
Quando $u(a+h) \neq u(a)$ para $h$ pequeno, escreva
\[
\frac{g(u(a+h)) - g(u(a))}{h}
= \frac{g(u(a+h)) - g(u(a))}{u(a+h) - u(a)} \cdot \frac{u(a+h) - u(a)}{h}.
\]
Quando $h \to 0$, $u(a+h) \to u(a)$, de modo que o primeiro fator tende a
$g'(u(a))$ e o segundo a $u'(a)$. (Uma demonstração completa precisa
tratar o caso em que $u(a+h) = u(a)$ para $h$ arbitrariamente pequeno;
esse ponto técnico é tratado na graduação.)
\end{proof}

\begin{example}
As \omterm{def:g12:deriv:derivative}{derivadas} usuais, válidas nos \omterm{def:g11:func:function}{domínios} naturais:
\[
(x^n)' = nx^{n-1}, \quad
\left(\frac{1}{x}\right)' = -\frac{1}{x^2}, \quad
(\sqrt{x})' = \frac{1}{2\sqrt{x}}, \quad
(\cos x)' = -\sin x, \quad
(\sin x)' = \cos x .
\]
A primeira demonstra-se por indução a partir da regra do produto, e as
duas últimas no \cref{ch:g12:trigo}.
\end{example}

\section{Variação e extremos}

\begin{theorem}[Sinal da derivada e variação]\label{thm:g12:deriv:variations}
Seja $f$ \omterm{def:g12:deriv:derivative}{derivável} em um \omterm{def:g10:numbers:interval}{intervalo} $I$.
\begin{enumerate}
  \item $f$ é \omterm{def:g12:seq:monotonic}{crescente} em $I$ se, e somente se, $f' \geq 0$ em $I$.
  \item $f$ é \omterm{def:g10:functions:variations}{decrescente} em $I$ se, e somente se, $f' \leq 0$ em $I$.
  \item Se $f' > 0$ em $I$, exceto em um número finito de pontos onde se
        anula, então $f$ é estritamente \omterm{def:g12:seq:monotonic}{crescente} em $I$.
\end{enumerate}
\end{theorem}

\begin{proof}[Demonstração parcial]
Se $f$ é \omterm{def:g12:seq:monotonic}{crescente}, todo quociente de diferenças
$\frac{f(a+h)-f(a)}{h}$ é $\geq 0$, e passar ao limite dá
$f'(a) \geq 0$. As implicações recíprocas apoiam-se no teorema do valor
médio, demonstrado na graduação; elas são \emph{admitidas neste nível}.
\end{proof}

\begin{definition}[Extremo local]\label{def:g12:deriv:extremum}
$f$ tem um \emph{máximo local}\index{extremo} em $a$ se $f(x) \leq f(a)$
para todo $x$ perto de $a$; os \omterm{def:g10:functions:extrema}{mínimos} locais definem-se de modo
simétrico.
\end{definition}

\begin{proposition}[Condição de primeira ordem]\label{prop:g12:deriv:critical}
Se $f$ é \omterm{def:g12:deriv:derivative}{derivável} em um \omterm{def:g10:numbers:interval}{intervalo} \emph{aberto} $I$ e tem um \omterm{def:g12:deriv:extremum}{extremo}
local em $a \in I$, então $f'(a) = 0$. A recíproca é falsa
($f(x) = x^3$ em $a = 0$).
\end{proposition}

\begin{proof}
Digamos que $a$ seja um \omterm{def:g12:deriv:extremum}{máximo local}. Para $h > 0$ pequeno,
$\frac{f(a+h)-f(a)}{h} \leq 0$, de modo que fazer $h \to 0^+$ dá
$f'(a) \leq 0$; para $h < 0$ o quociente é $\geq 0$, o que dá
$f'(a) \geq 0$. Logo $f'(a) = 0$.
\end{proof}

\begin{omfigure}
\begin{tikzpicture}
\begin{axis}[omaxis,
  width=8.4cm, height=5.6cm,
  xmin=-2.5, xmax=2.5, ymin=-4.2, ymax=4.2,
  xtick={-1,1}, ytick={-2,2}]
  \addplot[omProp, thick, domain=-1.8:-0.2] {2};
  \addplot[omProp, thick, domain=0.2:1.8] {-2};
  \addplot[omDef, thick, domain=-2.15:2.15] {x^3 - 3*x};
  \fill (axis cs:-1,2) circle (1.5pt);
  \fill (axis cs:1,-2) circle (1.5pt);
\end{axis}
\end{tikzpicture}

{\small Em um \omterm{def:g12:deriv:extremum}{extremo} local no interior de um \omterm{def:g10:numbers:interval}{intervalo} aberto, a
\omterm{def:g12:deriv:tangent}{tangente} (laranja) é horizontal: $f'(a) = 0$. Aqui $f(x) = x^3 - 3x$, com
\omterm{def:g12:deriv:extremum}{máximo local} em $-1$ e \omterm{def:g10:functions:extrema}{mínimo} local em $1$.}
\end{omfigure}

\begin{method}[Quadro de variação]\label{met:g12:deriv:table}
Para estudar uma \omterm{def:g11:func:function}{função} $f$: determine seu \omterm{def:g11:func:function}{domínio} e os limites na
fronteira; calcule $f'$ e fatore-a; determine o sinal de $f'$ em cada
subintervalo; registre tudo em um \omterm{def:g10:functions:table}{quadro de variação}, marcando os
\omterm{def:g12:deriv:extremum}{extremos}; deduza o número de soluções de $f(x) = k$ usando o teorema da
bijeção (\cref{thm:g12:limcont:bijection}).
\end{method}

\section{Convexidade}

\begin{definition}[Função convexa]\label{def:g12:deriv:convex}
Uma \omterm{def:g11:func:function}{função} $f$ definida em um \omterm{def:g10:numbers:interval}{intervalo} $I$ é
\emph{convexa}\index{convexidade} em $I$ se toda corda do seu \omterm{def:g10:functions:graph}{gráfico} fica
acima do \omterm{def:g10:functions:graph}{gráfico}: para todos $a, b \in I$ e $t \in \intcc{0}{1}$,
\[
f\bigl((1-t)a + t b\bigr) \leq (1-t) f(a) + t f(b).
\]
Ela é \emph{côncava}\index{côncava} se vale a desigualdade contrária
($-f$ convexa).
\end{definition}

\begin{omfigure}
\begin{tikzpicture}
\begin{axis}[omaxis,
  width=8.8cm, height=6cm,
  xmin=-1.9, xmax=2.5, ymin=-0.8, ymax=4.3,
  xtick={-1,1.5}, xticklabels={$a$,$b$}, ytick=\empty]
  \addplot[omThm, thick, domain=-1.25:1.75] {0.5*x + 1.5};
  \addplot[omProp, thick, domain=-0.45:2.3] {x - 0.25};
  \addplot[omDef, thick, domain=-1.6:2.0] {x^2};
  \draw[black, dashed, thin] (axis cs:0.25,0.0625) -- (axis cs:0.25,1.625);
  \fill (axis cs:-1,1) circle (1.5pt);
  \fill (axis cs:1.5,2.25) circle (1.5pt);
  \fill (axis cs:0.5,0.25) circle (1.5pt);
  \node[omThm, font=\small, above left] at (axis cs:1.7,2.4) {corda};
  \node[omProp, font=\small, below right] at (axis cs:1.75,1.5) {tangente};
\end{axis}
\end{tikzpicture}

{\small Uma \omterm{def:g11:func:function}{função} \omterm{def:g12:deriv:convex}{convexa}: toda corda (vermelha) fica acima do \omterm{def:g10:functions:graph}{gráfico}, e
o \omterm{def:g10:functions:graph}{gráfico} fica acima de cada uma de suas \omterm{def:g12:deriv:tangent}{tangentes} (laranja). O segmento
tracejado mostra a diferença entre $f\bigl((1-t)a+tb\bigr)$ e
$(1-t)f(a)+tf(b)$.}
\end{omfigure}

\begin{theorem}[Convexidade e derivadas]\label{thm:g12:deriv:convexity}
Seja $f$ \omterm{def:g12:deriv:derivative}{derivável} em um \omterm{def:g10:numbers:interval}{intervalo} $I$. São equivalentes:
\begin{enumerate}
  \item $f$ é \omterm{def:g12:deriv:convex}{convexa} em $I$;
  \item $f'$ é \omterm{def:g12:seq:monotonic}{crescente} em $I$;
  \item o \omterm{def:g10:functions:graph}{gráfico} de $f$ fica acima de cada uma de suas \omterm{def:g12:deriv:tangent}{tangentes}.
\end{enumerate}
Se $f$ é duas vezes \omterm{def:g12:deriv:derivative}{derivável}, isso também equivale a $f'' \geq 0$ em
$I$.
\end{theorem}

\begin{proof}[Demonstração de (2) $\Rightarrow$ (3)]
Fixe $a \in I$ e ponha $g(x) = f(x) - f(a) - f'(a)(x-a)$, a diferença
entre a curva e a \omterm{def:g12:deriv:tangent}{tangente} em $a$. Então $g$ é \omterm{def:g12:deriv:derivative}{derivável} e
$g'(x) = f'(x) - f'(a)$. Se $f'$ é \omterm{def:g12:seq:monotonic}{crescente}, $g' \leq 0$ em
$I \cap \intoc{-\infty}{a}$ e $g' \geq 0$ em
$I \cap \intco{a}{+\infty}$: $g$ decresce antes de $a$ e cresce depois,
de modo que $g$ atinge seu \omterm{def:g10:functions:extrema}{mínimo} $g(a) = 0$, logo $g \geq 0$. As demais
implicações são admitidas neste nível; a equivalência com $f'' \geq 0$
segue do \cref{thm:g12:deriv:variations} aplicado a $f'$.
\end{proof}

\begin{definition}[Ponto de inflexão]\label{def:g12:deriv:inflection}
Um ponto em que a curva de $f$ atravessa a sua \omterm{def:g12:deriv:tangent}{tangente} é um \emph{ponto
de inflexão}\index{ponto de inflexão}. Para $f$ duas vezes \omterm{def:g12:deriv:derivative}{derivável}, os
pontos de inflexão são os pontos em que $f''$ muda de sinal.
\end{definition}

\begin{example}\label{ex:g12:deriv:cube}
$f(x) = x^3$ tem $f''(x) = 6x$: $f$ é \omterm{def:g12:deriv:convex}{côncava} em $\intoc{-\infty}{0}$,
\omterm{def:g12:deriv:convex}{convexa} em $\intco{0}{+\infty}$, com \omterm{def:g12:deriv:inflection}{ponto de inflexão} na origem, onde a
curva atravessa sua \omterm{def:g12:deriv:tangent}{tangente} (o eixo $x$).
\end{example}

\begin{omfigure}
\begin{tikzpicture}
\begin{axis}[omaxis,
  width=7.6cm, height=5.6cm,
  xmin=-1.7, xmax=1.7, ymin=-2.9, ymax=2.9,
  xtick={-1,1}, ytick={-1,1}]
  \addplot[omProp, thick, domain=-1.2:1.2] {0};
  \addplot[omDef, thick, domain=-1.4:1.4] {x^3};
  \fill (axis cs:0,0) circle (1.5pt);
  \node[font=\small, anchor=east] at (axis cs:-0.45,-1.7) {côncava};
  \node[font=\small, anchor=west] at (axis cs:0.45,1.7) {convexa};
\end{axis}
\end{tikzpicture}

{\small O \omterm{def:g12:deriv:inflection}{ponto de inflexão} de $x \mapsto x^3$ na origem: a curva
atravessa sua \omterm{def:g12:deriv:tangent}{tangente} (laranja), \omterm{def:g12:deriv:convex}{côncava} à esquerda e \omterm{def:g12:deriv:convex}{convexa} à
direita.}
\end{omfigure}

\begin{example}[Desigualdades de convexidade]\label{ex:g12:deriv:ineq}
A exponencial é \omterm{def:g12:deriv:convex}{convexa} em $\R$ (sua segunda \omterm{def:g12:deriv:derivative}{derivada} é ela mesma,
positiva). Sua \omterm{def:g12:deriv:tangent}{tangente} em $0$ é $y = 1 + x$, de modo que
\[
\eu^x \geq 1 + x \quad \text{para todo } x \in \R .
\]
Essas \emph{desigualdades da reta \omterm{def:g12:deriv:tangent}{tangente}} são um produto padrão da
\omterm{def:g12:deriv:convex}{convexidade}.
\end{example}

\begin{method}[Usar a convexidade]\label{met:g12:deriv:convexity}
\begin{itemize}
  \item Para demonstrar uma desigualdade $f(x) \geq ax + b$: identifique a
        reta como \omterm{def:g12:deriv:tangent}{tangente} de uma $f$ \omterm{def:g12:deriv:convex}{convexa} e invoque o
        \cref{thm:g12:deriv:convexity}(3).
  \item Para localizar \omterm{def:g12:deriv:inflection}{pontos de inflexão}: resolva $f''(x) = 0$ \emph{e}
        verifique que $f''$ muda de sinal ali.
  \item Em um \omterm{def:g10:functions:table}{quadro de variação}, a \omterm{def:g12:deriv:convex}{convexidade} refina o esboço da curva
        (para que lado ela se curva).
\end{itemize}
\end{method}

\section{Exercícios}

\begin{exercise}[$\star$]\label{exo:g12:deriv:1}
Derive as funções seguintes em seus \omterm{def:g11:func:function}{domínios}:
\[
f(x) = x^3 - 5x + 2, \quad
g(x) = \frac{x}{x^2+1}, \quad
h(x) = (2x+1)^5, \quad
k(x) = \sqrt{x^2 + 1}.
\]
\end{exercise}

\begin{exercise}[$\star$]\label{exo:g12:deriv:2}
Dê a \omterm{def:g10:algebra:equation}{equação} da \omterm{def:g12:deriv:tangent}{tangente} à curva de $f(x) = x^2 - 3x + 1$ no ponto de
\omterm{def:g10:coordgeom:system}{abscissa} $a = 2$ e determine a posição da curva em relação a essa
\omterm{def:g12:deriv:tangent}{tangente}.
\end{exercise}

\begin{exercise}[$\star$]\label{exo:g12:deriv:3}
Usando a definição de \omterm{def:g12:deriv:derivative}{derivada}, calcule
\[
\lim_{x \to 0} \frac{(1+x)^{100} - 1}{x}
\qquad\text{e}\qquad
\lim_{h \to 0} \frac{\sqrt{4+h} - 2}{h}.
\]
\end{exercise}

\begin{exercise}[$\star\star$]\label{exo:g12:deriv:4}
Estude a \omterm{def:g11:func:function}{função} $f(x) = \dfrac{x^2 + 3}{x - 1}$ em seu \omterm{def:g11:func:function}{domínio}: limites,
\omterm{def:g12:deriv:derivative}{derivada}, \omterm{def:g10:functions:table}{quadro de variação}, \omterm{def:g12:deriv:extremum}{extremos} locais.
\end{exercise}

\begin{exercise}[$\star\star$]\label{exo:g12:deriv:5}
Uma caixa sem tampa é feita a partir de uma folha quadrada de papelão de
lado $30\,$cm, recortando-se quadrados iguais de lado $x$ nos cantos e
dobrando as laterais. Determine o $x$ que maximiza o volume da caixa.
\end{exercise}

\begin{exercise}[$\star\star$]\label{exo:g12:deriv:6}
Seja $f(x) = x^4 - 4x^3 + 10$.
\begin{enumerate}
  \item Calcule $f''$ e determine os \omterm{def:g10:numbers:interval}{intervalos} de \omterm{def:g12:deriv:convex}{convexidade} e os \omterm{def:g12:deriv:inflection}{pontos
        de inflexão} de $f$.
  \item Mostre que a \omterm{def:g12:deriv:tangent}{tangente} no \omterm{def:g12:deriv:inflection}{ponto de inflexão} de \omterm{def:g10:coordgeom:system}{abscissa} $2$
        atravessa a curva ali.
\end{enumerate}
\end{exercise}

\begin{exercise}[$\star\star$]\label{exo:g12:deriv:7}
Usando uma desigualdade da reta \omterm{def:g12:deriv:tangent}{tangente}, mostre que, para todo $x > 0$,
\[
\sqrt{x} \leq \frac{x + 1}{2},
\]
e identifique quando vale a igualdade. (Sugestão: a raiz quadrada é
\omterm{def:g12:deriv:convex}{côncava}.)
\end{exercise}

\begin{exercise}[$\star\star\star$]\label{exo:g12:deriv:8}
Seja $f$ \omterm{def:g12:deriv:convex}{convexa} e \omterm{def:g12:deriv:derivative}{derivável} em $\R$, e suponha que $f'$ se anule em algum
ponto $a$. Mostre que $f(a)$ é o \omterm{def:g10:functions:extrema}{mínimo} \emph{global} de $f$ em~$\R$.
\end{exercise}

\begin{exercise}[$\star\star\star$]\label{exo:g12:deriv:9}
Mostre que, entre todos os retângulos de perímetro fixo $p$, o quadrado
tem a maior área. Mostre em seguida que, entre todos os retângulos de área
fixa $A$, o quadrado tem o menor perímetro, e explique como as duas
afirmações se relacionam.
\end{exercise}

\section{Problema: O cálculo do salva-vidas}

\begin{problem}[{Problema de fim de semana --- o caminho de resgate
mais rápido obedece à lei de Snell da luz, a convexidade certifica
todo ótimo, e a lata de refrigerante ideal é exatamente tão alta
quanto larga}]
\label{pb:g12:deriv:1}
Uma salva-vidas avista um banhista em apuros --- na diagonal ao
longo da praia, dentro d'água. Correr é mais rápido que nadar: a
linha reta \emph{não} é o caminho mais rápido. O trajeto que
minimiza o tempo se dobra na beira d'água, e sua lei de refração é
precisamente a lei pela qual a luz se refrata ao entrar na água: a
natureza também deriva. Este problema treina a regra da cadeia
(\cref{thm:g12:deriv:chain}), executa o resgate, colhe as
desigualdades da \omterm{def:g12:deriv:convex}{convexidade} (\cref{thm:g12:deriv:convexity}) e
projeta uma lata.

\medskip
\noindent\textbf{Parte I --- Fluência com a regra da cadeia.}
\begin{enumerate}
  \item Derive: $\left(3x^2 + 1\right)^5$; \ $\sqrt{x^2 + 9}$; \
        $\dfrac{1}{x^2 + 1}$.
  \item Dê a \omterm{def:g12:deriv:tangent}{tangente} a $y = \sqrt{x^2 + 9}$ em $x = 4$.
  \item Monte o \omterm{def:g10:functions:table}{quadro de variação} de $f(x) = \dfrac{x}{x^2 + 1}$
        em $\R$ (\omterm{def:g12:deriv:extremum}{extremos} incluídos).
  \item Para $g(x) = x^3 - 3x^2 + 4$: \omterm{def:g10:functions:table}{quadro de variação} \emph{e}
        quadro de \omterm{def:g12:deriv:convex}{convexidade} ($g''$), com o \omterm{def:g12:deriv:inflection}{ponto de inflexão}
        (\cref{def:g12:deriv:inflection}).
  \item Calcule a \omterm{def:g12:deriv:tangent}{tangente} a $g$ em seu \omterm{def:g12:deriv:inflection}{ponto de inflexão} e
        mostre que $g(x) - (\text{tangente}) = (x - 1)^3$: o que
        a mudança de sinal diz sobre o modo como a \omterm{def:g12:deriv:tangent}{tangente}
        encontra a curva em uma inflexão?
\end{enumerate}

\medskip
\noindent\textbf{Parte II --- O resgate.}
Beira d'água ao longo do eixo $x$; a salva-vidas está na areia, em
$A(0, 30)$, e o banhista espera em $B(40, -20)$ (em metros). Ela
corre a $5$ m/s na areia, nada a $2$ m/s e entra na água em um
ponto $(x, 0)$ à sua escolha.
\begin{enumerate}[resume]
  \item Modelagem: exprima o tempo total de resgate $T(x)$ e diga
        por que só $x \in \intcc{0}{40}$ merece consideração.
  \item Derive $T$ (a regra da cadeia em ação --- duas vezes).
  \item Seja $\theta_1$ o ângulo do trecho de corrida com a
        \emph{perpendicular} à beira d'água, e $\theta_2$ o do
        trecho de natação. Mostre que
        $\sin\theta_1 = \frac{x}{\sqrt{x^2 + 900}}$ e que a
        condição $T'(x) = 0$ se lê
        \[
        \frac{\sin\theta_1}{5} = \frac{\sin\theta_2}{2}
        \]
        --- a \emph{lei de Snell}, com as duas velocidades no
        lugar das velocidades da luz no ar e na água.
  \item O segmento reto $[AB]$ cruza a beira d'água em $x = 24$.
        Calcule $T(24)$, $T(40)$ (correr pela praia e depois
        nadar em linha reta) e $T(0)$: a reta geométrica é a mais
        rápida? Alguma das estratégias extremas é?
  \item Localize o ponto ótimo de entrada por \omterm{thm:g12:limcont:ivt}{dicotomia} em $T'$
        (ele é $\approx 33.7$ m): dê $x^*$ com precisão de metros
        e o tempo recorde $T(x^*)$ com precisão de décimos de
        segundo. Quanto a curvatura ótima economiza em relação à
        linha reta?
  \item Por que o ponto crítico é um \omterm{def:g10:functions:extrema}{mínimo} \emph{global}? (Cada
        parcela de $T$ é \omterm{def:g12:deriv:convex}{convexa} --- admita que uma soma de
        funções \omterm{def:g12:deriv:convex}{convexas} é \omterm{def:g12:deriv:convex}{convexa} --- e aplique o
        \cref{exo:g12:deriv:8}.)
  \item O princípio de Fermat afirma que a luz percorre sempre o
        caminho de menor \emph{tempo}. Deduza por que um raio de
        luz se dobra exatamente na superfície ao entrar na água
        (onde a luz é mais lenta) e nomeie a observação cotidiana
        que isso explica sobre um lápis dentro de um copo com
        água.
\end{enumerate}

\medskip
\noindent\textbf{Parte III --- A colheita da \omterm{def:g12:deriv:convex}{convexidade}.}
\begin{enumerate}[resume]
  \item Mostre que $x \mapsto x^2$ é \omterm{def:g12:deriv:convex}{convexa} e demonstre duas
        vezes a desigualdade do \omterm{prop:g10:coordgeom:midpoint}{ponto médio}
        $\left(\frac{a + b}{2}\right)^2 \leq
        \frac{a^2 + b^2}{2}$: uma vez pela \omterm{def:g12:deriv:convex}{convexidade}, outra
        expandindo $(a - b)^2 \geq 0$.
  \item Deduza que a \emph{\omterm{def:g11:stat:mean}{média} quadrática}
        $\sqrt{\frac{a^2 + b^2}{2}}$ domina a \omterm{def:g11:stat:mean}{média} aritmética e
        monte a cadeia completa da série --- harmônica $\leq$
        geométrica $\leq$ aritmética $\leq$ quadrática ---,
        verificando as quatro em $a = 2$, $b = 8$.
  \item \omterm{def:g12:deriv:tangent}{Tangente} por baixo da curva: $x \mapsto \frac1x$ é
        \omterm{def:g12:deriv:convex}{convexa} em $\intoo{0}{+\infty}$; escreva sua \omterm{def:g12:deriv:tangent}{tangente} em
        $1$ e deduza a desigualdade $\frac1x \geq 2 - x$ para
        todo $x > 0$. Onde ela é igualdade?
  \item Inflexão em campo: durante uma epidemia, o número
        acumulado de casos segue uma curva em S. O que acontece
        no seu \omterm{def:g12:deriv:inflection}{ponto de inflexão}, e por que essa data é a que os
        epidemiologistas aguardam? Relacione com o sinal da
        segunda \omterm{def:g12:deriv:derivative}{derivada} de cada lado.
\end{enumerate}

\medskip
\noindent\textbf{Parte IV --- A lata ideal.}
Uma lata cilíndrica deve conter $V = 330$ cm$^3$ com o \omterm{def:g10:functions:extrema}{mínimo} de
metal: superfície $S = 2\pi r^2 + 2\pi r h$, com a restrição
$\pi r^2 h = V$.
\begin{enumerate}[resume]
  \item Exprima $S(r) = 2\pi r^2 + \dfrac{2V}{r}$, derive e
        calcule o raio e a altura ótimos para $V = 330$ (com
        precisão de milímetros).
  \item Demonstre a elegante lei geral: no ótimo, $h = 2r$ --- a
        lata ideal é exatamente tão alta quanto larga.
  \item Verifique que $S''(r) > 0$ e conclua (de novo pelo
        \cref{exo:g12:deriv:8}) que o ótimo é global. As latas de
        refrigerante reais são visivelmente mais altas que
        largas: aponte uma razão não matemática.
  \item Final --- a linha de montagem da otimização: modelar,
        perguntar à \omterm{def:g12:deriv:derivative}{derivada}, resolver os pontos críticos,
        certificar com \omterm{def:g12:deriv:convex}{convexidade} ou com um quadro, interpretar.
        Percorra a lista para a salva-vidas, para a lata e para a
        viga mais forte do ano 11 (\cref{pb:g11:deriv:1}) --- e
        diga o que o princípio de Fermat afirma sobre quem mais
        percorre essa mesma linha de montagem.
\end{enumerate}
\end{problem}
